\chapter{第三讲\quad 广义相对论与拟线性波动方程}

我们从研究拟线性波动方程最优适定性问题的主要动机开始.
\section*{3.1 广义相对论中的柯西问题}

爱因斯坦广义相对论的主要对象是时空. 它可定义为可微、定向的四维 Lorentz 流形 $(\mathcal{M}, g)$ 的等价类. 称两个 Lorentz 流形 $(\mathcal{M}, g)$、$(\mathcal{M}', g')$ 等价, 若存在微分同胚 $\mathcal{M} \longrightarrow \mathcal{M}'$ 使 $g' = \Phi_* g$. 时空就是此类 Lorentz 流形的一个等价类. 时空度量 $g$ 须满足爱因斯坦场方程
\[
R_{\alpha\beta}-\frac{1}{2}g_{\alpha\beta}R=T_{\alpha\beta},
\]
其中 $R_{\alpha\beta}$ 为 Ricci 曲率, $R$ 为度量的标量曲率, $T_{\alpha\beta}$ 为定义在 $(\mathcal{M}, g)$ 上某物质场的能量--动量张量. 我们限于真空特例, 即 $T\equiv0$, 此时方程变为
\[
R_{\alpha\beta}=0.
\]

爱因斯坦方程的演化特征比其他方程微妙得多, 因为并不存在时间变量 $t$ 的内在定义. 因此谈论初值问题似乎同样含混. 不过它可如下定义:
\setcounter{mtthm}{1}
\begin{definition}\label{def:3.2}
初数据集是一个三元组 $(\Sigma, \bar{g}, \bar{k})$, 由一个三维完备黎曼流形 $(\Sigma, \bar{g})$ 与一个在 $\Sigma$ 上的 2-协变对称张量 $k$ 组成, 满足约束方程:
\[
\bar{\nabla}^{j}\bar{k}_{i j}-\bar{\nabla}_{i}\mathrm{tr}\,\bar{k}=0,\qquad
\bar{R}-|\bar{k}|^{2}+(\mathrm{tr}\,\bar{k})^{2}=0,
\]
其中 $\bar{\nabla}$ 为协变导数, $\bar{R}$ 为 $(\Sigma, \bar{g})$ 的标量曲率. 标量 $\bar{k}$ 为 $\bar{k}$ 相对 $\bar{g}$ 的迹.\end{definition}

\setcounter{mtthm}{2}
\begin{definition}\label{def:3.3}
若初数据集对应于 Minkowski 空间中具诱导度量与第 II 基本形式的完备类空超曲面, 则称其为平坦(或平凡)的. 若在某坐标系 $(x^1, x^2, x^3)$ 定义于 $\Sigma$ 上无穷远邻域$^{13}$ 中, 使得度量 $\bar{g}$ 趋于欧氏度量、$\bar{k}$ 趋于零$^{14}$, 则称初数据集为渐近平坦(AF):
\[
r=\sqrt{(x^{1})^{2}+(x^{2})^{2}+(x_{3})^{2}}\to\infty.
\]
\end{definition}

\setcounter{mtthm}{3}
\begin{definition}\label{def:3.4}
初数据集 $(\Sigma, \bar{g}, \bar{k})$ 的柯西发展是一个时空流形 $(\mathcal{M}, g)$, 满足爱因斯坦方程以及一个嵌入$^{15}$ $i: \Sigma \longrightarrow \mathcal{M}$, 使得 $\bar{g}, \bar{k}$ 为 $\Sigma$ 在 $\mathcal{M}$ 中的诱导第 I、第 II 基本形式. 为保证解对数据的唯一依赖, 还要求该发展是整体双曲的$^{16}$.\end{definition}

爱因斯坦方程演化特征的微妙性, 使得对初值问题的局部存在唯一性(对很大一类演化方程已满意解决)成为首要问题. 对爱因斯坦方程, 此类结果最早由 Y. Choquet-Bruhat \cite{klai7} 借助波坐标得到, 波坐标使她能把爱因斯坦方程写成非线性波动方程组的形式, 从而可应用只基于能量估计与 $L^\infty \subset H^s$ Sobolev 不等式的标准局部存在理论.
\setcounter{mtthm}{4}
\begin{definition}\label{def:3.5}
若坐标 $x^\alpha$ 满足 $\Box_g x^\alpha = 0$, 则称其为波坐标. 给定这样的坐标系, 爱因斯坦真空方程取形式
\[
g^{\alpha\beta}\partial_{\alpha}\partial_{\beta}g_{\mu\nu}=N_{\mu\nu}(g_{\alpha\beta},\partial g_{\alpha\beta}).
\]
\end{definition}

对非线性微分方程组, 局部存在唯一性结果通过简单的延拓论证导出整体结果. 对爱因斯坦方程表述同类结果稍显微妙; 这由 Y. Choquet-Bruhat 与 R.P. Geroch 在 \cite{klai8} 中完成. 该结果断言: 对每个初数据集存在唯一的极大未来发展. 因此, 由演化方法从特定初数据集出发构造的任何发展, 必包含于其极大发展中. 这可以说解决了广义相对论中整体$^{17}$ 存在唯一性问题; 所有进一步问题可说关乎极大发展的定性性质. 核心问题是奇点的存在性与特征. 首先, 我们可定义正则极大发展为完备者, 即所有未来类时与零测地线都完备. Penrose 的一个著名定理断言, 在初数据集的某些相当极端的条件下$^{18}$, 相应的未来极大发展必不完备. 此外, 除平坦 Minkowski 空间本身外, 爱因斯坦真空方程的所有显式 AF 解(如 Kerr 族)都有奇点. 与 Penrose 定理相对的是如下完备性结果, 见 \cite{klai12}.
\setcounter{mtthm}{5}
\begin{theorem}[Minkowski 空间的全局稳定性]\label{thm:3.6}
任何充分接近平凡数据的渐近平坦初数据集, 具有完备的极大未来发展 $(M,g)$. 而且时空 $M$ 在如下意义下接近 Minkowski 空间: 其曲率张量沿任何类时或零测地线趋于零. 然而其因果结构渐近非平凡.\end{theorem}

\begin{problem*}[闵可夫斯基空间的强稳定性问题]
是否存在一个尺度不变的小性条件$^{20}$, 使得所有初数据集 $(\Sigma, g, k)$ 满足它的发展都有完备的极大未来发展? 
\end{problem*}

\setcounter{mtthm}{6}
\begin{remark}\label{rem:3.7}
\thmref{thm:3.6} 的旧稳定性结果与强稳定性之间的关系, 可与 Tao 的结果 \cite{klai71} 之后波动映射的情形相比. 事实上 Sideris \cite{klai63} 的一个较早结果对含大量导数的加权 Sobolev 范数中的小数据给出整体存在. Tao 的新结果只要求含初值映射最少导数的尺度不变范数的小性.\end{remark}

\setcounter{mtthm}{7}
\begin{remark}\label{rem:3.8}
Minkowski 时空的强稳定性本身应被视为通向证明著名的 Penrose 宇宙监督猜想的更高目标的中间步骤. 我们在第三讲的补充中陈述这些猜想.\end{remark}

证明 Minkowski 空间的强稳定性是一个长期目标. 下面我们提出一个远为简单、但仍极具挑战性的中间猜想, 它启发下一节的主题.
\begin{theorem*}[有限 $L^{2}$ 曲率猜想]
Bruhat--Geroch 结果可推广到初数据集 $(\Sigma, g, k)$ 满足
\[
R(g)\in L^{2}_{(loc)}(\Sigma)\qquad \text{且}\qquad k\in H^{1}_{(loc)}(\Sigma).
\]
\end{theorem*}

$L^2$ 有界性猜想是我们此刻就 Minkowski 空间强稳定性猜想所能期望的最合理目标. 它可与 $\mathbb{R}^{1+3}$ 中 YM 方程局部(在此情形也是整体的! )适定性的证明相比, 见 \cite{klai34}. 但它要复杂得多. 我们或许先满足于一个更简单的目标: 在 $H^{2+\epsilon}(\mathbb{R}^3)$ 中证明系统 (64) 的适定性. 为稍加简化, 考虑形式
\[
G^{\alpha\beta}(\phi)\partial_{\alpha}\partial_{\beta}\phi=N(\phi,\partial\phi)
\]
的波动方程, 满足初值条件 $\phi[0] \in H^s$. 这里 $G^{\alpha\beta}(\phi)$ 是 $\mathbb{R}^{1+3}$ 中一族 Lorentz 度量, 在原点邻域光滑依赖于 $\phi$, 而 $N$ 在 $\partial\phi$ 中为二次. 我们想证明对 $s > 2 + \epsilon$ 该 IVP 适定. 对形如
\[
\Box\phi=N(\phi,\partial\phi)
\]
的半线性波动方程的类似结果极易证明$^{21}$, 它只依赖于如下形式的 Strichartz 型估计:
\[
\|\partial\psi\|_{L_{t}^{2}L_{x}^{\infty}([0,T]\times\mathbb{R}^{n})}\lesssim\|\phi[0]\|_{H^{s}(\mathbb{R}^{n})}
\]
对任意 $s > \frac{n}{2} + \frac{1}{2}$、$n = 3$ 成立. 特别地它不要求非线性项满足如零条件的任何特殊结构.
\section*{3.9 拟线性波动方程的最优适定性}

我打算报告我与 I. Rodnianski \cite{klai42} 关于拟线性波动方程最优适定性问题的近期结果. 受形如 (64) 与 (67) 的拟线性波动方程启发, 我们考虑 $\mathbb{R}^{n+1}$($n \geq 3$)中的简化模型方程
\[
\partial_{t}^{2}\phi-g^{i j}(\phi)\partial_{i}\partial_{j}\phi=N(\phi,\partial\phi),
\]
其中 $N$ 为在 $\partial\phi$ 中二次的光滑函数. 基于 $H^{s}$ 能量估计与 Sobolev 嵌入定理的经典局部存在唯一性结果表明, (67) 在 Sobolev 空间 $H^{s}$(任意 $s > \frac{n}{2} + 1$)中局部适定. 由标准分部积分与交换子引理得到的 $H^{s}$ 能量估计取形式
\[
\|\phi(t)\|_{H^{s}}\leq M(\|\phi_{0}\|_{H^{s}}+\|\phi_{1}\|_{H^{s-1}}),
\]
其中 $M$ 连续依赖于积分 $\int_0^t \|\partial\phi\|_{L_x^\infty}$. 对充分小的时刻 $t$, Sobolev 估计 $\|\partial\phi\|_{L^\infty} \leq c\|\phi\|_{H^s}$($s > \frac{n}{2} + 1$)使我们能把上式右端纯用 $\|\phi_0\|_{H^s} + \|\phi_1\|_{H^{s-1}}$ 界定. 这些对 $s > \frac{n}{2} + 1$ 的局部先验估计可轻易化为局部存在唯一性结果.
量 $\int_0^t \|\partial\phi\|_{L_\infty^\infty}$ 可借助 Strichartz 估计 (66) 更好地控制. 这就带来了上述 $\frac{1}{2}$ 导数的获得.
在拟线性情形, 度量 $g^{ij}$ 依赖于解 $\phi$, 因而其正则性至多与 $\phi$ 相同. 因此需处理对线性方程 $\Box_g\phi = \partial_t^2 - g^{ij}\partial_t\partial_j\phi = 0$ 的 Strichartz 估计的证明, 其中系数与线性化波动算子的系数一致. 线性方程 $\Box_g$ 的系数所需条件须与 (67) 中非线性系数 $g^{ij}(\phi)$ 对 $\phi$ 的依赖一致. 鉴于 $\|\phi\|_{H^s}$ 的 Sobolev 范数应有界, 我们可设$^{22}$ $\|g\|_{H^s}$ 有界. 而且由于拟线性问题 (67) 在时刻 $t$ 的解的 $H^s$ 范数被初数据的 $H^s$ 范数控制, 只要 $\int_0^t\|\partial\phi\|_{L_\infty^\infty}$ 有界, 我们可归纳地假设度量 $g^{ij} = g^{ij}(\phi)$ 也满足 $\int_0^t\|\partial g\|_{L_\infty^\infty} \leq B_0$. 依我们在平坦空间 Strichartz 估计的经验, 我们不期望直接触及 $\|\partial\phi\|_{L_t^1L_\infty^\infty}$ 的范数, 但在局部时间它可由 $\|\partial\phi\|_{L_t^2L_\infty^\infty}$ 控制. 因此, 为闭合论证并导出 (67) 改进的局部存在唯一性结果, 我们需要证明 $H^s$ 范数的有界性, 以及对线性方程 $\Box_g\phi = 0$ 解的局部 $L_t^2L_\infty^\infty$-Strichartz 估计.
历史注记: 关于对系数粗糙$^{23}$ 的 $\Box_g \phi = 0$ 解之 Strichartz 型估计的首要重要工作归于 H. Smith \cite{klai60}. 他证明了标准 Strichartz 估计对 $C^{1,1}$ 系数方程成立. 但 $C^{1,1}$ 条件对应用于拟线性方程过强. 而且 H. Smith 与 C. Sogge \cite{klai61} 的一些重要反例表明, 若系数比 $C^{1,1}$ 更粗糙, 标准 Strichartz 估计失效.
这是 Bahouri 与 J.-Y. Chemin 的重要突破之前的状况. 在 \cite{klai1} 中他们证明, 对某些比 $C^{1,1}$ 更粗糙、与拟线性方程应用相容的系数, Strichartz 估计的较弱形式仍然存在. 即他们为变系数波动方程 $\Box_g \psi = 0$(系数满足上述类型假设)建立了带损失的 Strichartz 估计
\[
\left\|\partial\psi\right\|_{L_{t}^{2}L_{x}^{\infty}}\lesssim\left\|\psi[0]\right\|_{H^{\frac{n}{2}+\frac{1}{2}+\sigma}},
\]
只要损失 $\sigma \leq \frac{1}{2}$, 此估计可用于证明此处拟线性问题的局部适定性结果, 强于经典结果. H. Bahouri 与 J.-Y. Chemin 的结果基于对线性化问题 $\Box_g \psi = 0$(度量 $g$ 满足上述类型条件)带损失 $\sigma > \frac{1}{4}$ 的 Strichartz 估计. 后来 D. Tataru \cite{klai72} 证明 Strichartz 估计中的损失$^{24}$ 可改进为 $\sigma > \frac{1}{6}$. 事实上 \cite{klai72} 给出了低于 $C^{1,1}$ 时度量的光滑性与 Strichartz 估计中相应损失之间的精确关系. 最近 H. Smith 与 D. Tataru \cite{klai62} 证明这些结果也是尖锐的, 因此仅基于线性理论无法改进 Tataru 的结果.
Bahouri--Chemin 与 Tataru 的结果都基于两个主要要素.
a) 仿微分法: 考虑对应于空间频率 $\frac{\lambda}{2} \leq |\xi| \leq 2\lambda$($\lambda \in 2^\mathbf{N}$)截断的二进投影算子 $P_\lambda$. 令 $\phi^\lambda = \phi_\lambda = P_\lambda \phi$, 有 $\phi = \sum_\lambda \phi_\lambda$. 若 $\phi$ 满足拟线性波动方程 (67), 则 $\phi^\lambda$ 满足方程
\[
\Box_{g_{\leq\lambda^{a}}}\phi^{\lambda}=-\partial_{t}^{2}\phi^{\lambda}+g_{\leq\lambda^{a}}^{i j}\partial_{i}\partial_{j}\phi^{\lambda}=R_{\lambda}^{a},
\]
其中 $g_{\leq\lambda^a} = \sum_{\mu\leq\lambda^a} P_\mu g$, $0 < a < 1$, $R_\lambda$ 为易控制的误差项. Bahouri--Chemin 后来由 Tataru 精炼的原始思想, 是在频率依赖的时间区间 $I_{\lambda}$(大小 $|I_{\lambda}| \lesssim \lambda^{1-a}$)上, 对 (69) 证明无损失的 Strichartz 估计. 正是控制时间区间 $I_{\lambda}$ 大小上的这一损失, 导致了 (71) 中可微性的损失.
b) 基于参数化的色散不等式证明: 如我们所知, Strichartz 不等式证明的主要困难构件是能量与色散不等式. Bahouri--Chemin 与 Tataru 都利用参数化构造$^{25}$ 给出近似线性化方程 (69) 的解. 他们借助驻相法, 由参数化构造给出的近似解显式形式导出色散不等式.
c) 程函方程(eikonal equation): 两种方法的主要要素之一是构造程函方程的解
\[
u_{t}^{2}-g_{\leq\lambda^{a}}^{i j}\partial_{i}u\partial_{j}u=0,\qquad u(0,x,\xi)=x\cdot\xi,
\]
以及相伴的输运方程.
在 \cite{klai30} 中我发展了另一种更几何的方法, 基于第一讲讨论的向量场法的修改来证明 (69) 的色散不等式. 回忆命题 (1.9) 可知, 我们无需解的显式表示, 只用了标准波动方程与伸缩、推进及角动量向量场的交换性质, 就能证明色散不等式. 为对 (69) 型拟线性方程实施此策略, 我们需要类似的一族 Killing 与保形 Killing 向量场, 但对应度量 $g \leq \lambda^a$. 这当然不可能; 一般度量没有任何对称性. 我们能做的最好是构造形变张量尽可能小的向量场. 在 \cite{klai30} 中我概述了如何基于同一程函方程 (79) 但具不同初值(对应于此前的一族零超平面, 而此处对应于时间轴上顶点的一族出特征锥)来作此构造. 这一直接几何方法使我恢复了三维拟线性波动方程的 Tataru 结果.
利用 \cite{klai30} 中类似的几何方法, 最近 Rodnianski 与我在 \cite{klai42} 中, 通过考虑方程 (67) 的非线性特征, 改进了 Bahouri--Chemin 与 Tataru 的适定性结果. 我们不只对满足 $\|g\|_H^s$ 与 $\|\partial g\|_{L_1^1 L_\infty^s}$ 估计的 $\Box_g \phi = 0$ 解证明 Strichartz 估计, 而是本质地利用了方程 (67) 的系数 $g^{ij}(\phi)$ 自身满足形如 $\Box_g g_{ij} = N$(其中 $N$ 只依赖于 $\phi$ 与 $\partial \phi$)的方程这一事实.
我们的主要结果包含在如下定理中.
\setcounter{mtthm}{9}
\begin{theorem}\label{thm:3.10}
$\mathbb{R}^{3+1}$ 中的拟线性初值问题 (67) 在 $H^{s_*}$ 中局部适定, 其中 $s_* > s_0 = 2 + \frac{2 - \sqrt{3}}{2}$. 即对任意初值 $\phi[0] \in H^{s_*}$, 存在充分小的区间 $[0, T]$($T = T(\|\phi[0]\|_{H^{s_*}})$), 使得问题 (67) 有唯一解 $\phi \in C([0, T], H^{s_*})\cap C^{1}([0, T], H^{s_0-1})$. 此外 $\phi$ 满足 Strichartz 型估计
\[
\|\partial\phi\|_{L_{[0,T]}^{2}L_{x}^{\infty}}\leq c\, T^{s_{*}-s_{0}}\|\phi[0]\|_{H^{s_{*}}}.
\]
而且我们相信, 对 \cite{klai42} 的证明作相对较小修改, 我们现在已具备证明爱因斯坦方程$^{26}$ 在 $H^{2+\epsilon}$ 中局部适定所需的全部工具.\end{theorem}

\section*{3.11 \thmref{thm:3.10} 证明概要}

经过一系列如今已典范的简化(见本讲补充), 我们可把定理的证明化为如下情形:
设给定一族$^{27}$ 黎曼度量 $h_{ij}^{\lambda}(t,x)dx^{i}dx^{j}$, 定义在时间板 $I_{\lambda} \times \mathbb{R}^{3}$ 上, $I_{\lambda} = [0, t_{*}]$, $t_{*} \leq \lambda^{a}$, 满足如下条件
\[
\|\partial^{1+m}h_{\lambda}\|_{L_{I_{\lambda}}^{1}L_{x}^{\infty}}\lesssim\lambda^{-(1-a)(m+1)},\qquad
\|\partial^{1+m}h_{\lambda}\|_{L_{I_{\lambda}}^{2}L_{x}^{\infty}}\lesssim\lambda^{-\frac{2-a}{2}-(1-a)m},
\]
\[
\|\partial^{1+m}h_{\lambda}\|_{L_{I_{\lambda}}^{\infty}L_{x}^{\infty}}\lesssim\lambda^{-1+\frac{a^{2}}{2}-(1-a)m},\qquad
\|\partial_{x}^{\frac{1}{2}+m}(\partial^{2}h_{\lambda})\|_{L_{I_{\lambda}}^{\infty}L_{x}^{2}}\lesssim\lambda^{-1+\frac{a^{2}}{2}-(1-a)m},
\]
\[
\|\partial^{m}\bar{\Box}_{h_{\lambda}}h_{\lambda}\|_{L_{I_{\lambda}}^{1}L_{x}^{\infty}}\lesssim\lambda^{-(2-a)-(1-a)m},
\]
其中 $\bar{\Box}_{h_{\lambda}} = \partial_{t}^{2} - h_{\lambda}^{ij} \partial_{i} \partial_{j}$. 我们需证明对协变波动方程
\[
\Box_{h_{\lambda}}\psi=0
\]
的解的如下衰减估计, 其中
\[
\Box_{h_{\lambda}}\psi=\frac{1}{\sqrt{\det h_{\lambda}}}\partial_{t}\sqrt{\det h_{\lambda}}\partial_{t}\psi-\frac{1}{\sqrt{\det h_{\lambda}}}\partial_{i}\bigl(h_{\lambda}^{i j}\sqrt{\det h_{\lambda}}\partial_{j}\psi\bigr).
\]

\setcounter{mtthm}{11}
\begin{theorem}\label{thm:3.12}
设 $\psi$ 为 $[0, t_*] \times \mathbb{R}^n$ 中 (77) 的解, $t_* \leq \lambda^a$, 其初数据 $\psi[0]$ 支在物理空间中以原点为中心、半径 $\frac{1}{2}$ 的球内. 设度量 $h_\lambda$ 满足 (138)--(142) 且 $a < -1 + \sqrt{3}$. 则对所有 $t \in [0, t_*]$,
\[
\|\partial\psi(t)\|_{L^{\infty}}\leq(1+t)^{-1+\epsilon}\|\psi[0]\|_{H^{k}}.
\]
\end{theorem}

证明基于第三讲补充中讨论的\propref{prop:4.11} 的弯曲时空类似. 在陈述它之前我们需要一些定义:
\setcounter{mtthm}{12}
\begin{definition}\label{def:3.13}
零对(null pair)由两向量场 $L,\underline{L}$ 组成, 满足
\[
H(L,L)=H(\underline{L},\underline{L})=0,\qquad H(L,\underline{L})=-2.
\]
与零对 $L, \underline{L}$ 相伴的零标架(null frame)由 4 个线性无关向量 $e_1, e_2, e_3 = \underline{L}, e_4 = L$ 组成, 满足
\[
H(e_{A},L)=H(e_{A},\underline{L})=0,\quad H(e_{A},e_{B})=\delta_{A B},\quad\text{对 }A=1,2.
\]
\end{definition}

\setcounter{mtthm}{13}
\begin{definition}\label{def:3.14}
光学函数$^{28}$ $u$ 定义为程函方程
\[
H^{\alpha\beta}\partial_{\alpha}u\partial_{\beta}u=0,\qquad\text{或}\quad\partial_{t}u=|\nabla u|_{h}
\]
的向前解, 边界条件为在时间轴上 $u = t$. 时间函数 $t$ 的等值超曲面 $\Sigma_t$ 与 $u$ 的等值超曲面 $C_u$ 的交的 2-曲面记为 $S_{t,u}$.\end{definition}

\setcounter{mtthm}{14}
\begin{definition}\label{def:3.15}
给定这样的 $u$, 我们定义典范零对如下:
\[
L=T+N,\qquad \underline{L}=T-N,
\]
其中
\[
T=\partial_{t},\qquad N=-\frac{h^{i j}\partial_{j}u}{|\nabla u|_{h}}\partial_{i}.
\]
注意 $T + N$ 是零锥 $C_u$ 的生成法向量. 我们可补充 $L, \underline{L}$ 构成张成整个切空间的零标架. 这可通过在 $S_{t,u} = C_u \cap \Sigma_t$ 上任取一组正交标架 $e_A$ 实现. 我们重写完整零标架 $e_1, e_2, e_3 = \underline{L}, e_4 = L$.
我们还引入如下:
\begin{enumerate}[label=\roman*)]
\item 能量--动量张量: $Q_{\alpha\beta} = \partial_{\alpha}\psi\partial_{\beta}\psi - \tfrac{1}{2}H_{\alpha\beta}(H^{\mu\nu}\partial_{\mu}\psi\partial_{\nu}\psi)$.
\item 修正能量--动量张量: $\bar{Q}_{\alpha\beta} = Q_{\alpha\beta} + (n-1)t\psi\partial_{\beta}\psi - \frac{n-1}{2}\psi^{2}$.
\item 修正 Morawetz 向量场: 令 $K = \tfrac{1}{2}(\underline{u}^2L + u^2\underline{L})$, 其中 $\underline{u} = 2t - u$.
\item 修正形变张量: $\bar{\pi} = (K)\pi - 4tH$.
\item 共形能量: $Q_0[\psi](t) = \int_{\Sigma_t} \bar{Q}(K, \partial_t)[\psi]$.
\item 全共形能量: $Q[\psi](t) = Q_0[\psi](t) + Q_0[\partial_t\psi](t) + Q_0[\partial_t^2\psi](t)$.
\item $Q$ 的零分量: 相对零标架 $Q$ 的分量为
\[
Q_{44}=(D_{4}\psi)^{2},\quad Q_{34}=|\bar{\nabla}\psi|^{2},\quad
Q_{33}=(D_{3}\psi)^{2},\quad Q_{3A}=D_{3}\psi\bar{\nabla}_{A}\psi,
\]
\[
Q_{4A}=D_{4}\psi\bar{\nabla}_{A}\psi,\qquad Q_{A B}=\bar{\nabla}_{A}\psi\bar{\nabla}_{B}\psi+\frac{1}{2}(D_{3}\psi D_{4}\psi-|\bar{\nabla}\psi|^{2})\delta_{A B}.
\]
\item 共形能量密度: 由 vii) 易计算:
\[
Q(K,\partial_{t})[\psi]=\frac{1}{4}\bigl(\underline{u}^{2}(D_{4}\psi)^{2}+u^{2}(D_{3}\psi)^{2}+(\underline{u}^{2}+u^{2})|\nabla\psi|^{2}\bigr).
\]
\end{enumerate}
\end{definition}

\thmref{thm:3.12} 的证明可化为如下:
\setcounter{mtthm}{15}
\begin{theorem}[有界性定理]\label{thm:3.16}
在\thmref{thm:3.12} 同样假设下, 有
\[
\mathcal{Q}(t)\lesssim\mathcal{Q}(0).
\]
\end{theorem}

下面我们勾画有界性定理证明的一些主要思想. 我们限于讨论 $Q_{0}$ 的有界性.
步骤 1. 引入量
\[
\mathcal{E}_{0}[\psi](t)=\int_{\Sigma_{t}}\bigl(\underline{u}^{2}(D_{4}\psi)^{2}+u^{2}(D_{3}\psi)^{2}+\underline{u}^{2}|\bar{\nabla}\psi|^{2}+\psi^{2}\bigr)\lesssim\mathcal{Q}_{0}[\psi](t).
\]
由类似于\propref{prop:4.11} 之 iii) 的比较定理, 我们证明
\[
\mathcal{E}_{0}[\psi](t)\lesssim\mathcal{Q}_{0}[\psi](t).
\]

步骤 2. \thmref{thm:3.16} 证明的关键要素是如下\propref{prop:4.11} 的广义共形恒等式的弯曲时空版本.
\setcounter{mtthm}{16}
\begin{proposition}[广义共形恒等式]\label{prop:3.17}
设 $\psi$ 为 $\Box_h \psi = 0$ 的解. 则
\[
\mathcal{Q}_{0}[\psi](t)=\mathcal{Q}_{0}[\psi](t_{0})-\frac{1}{2}\int_{[t_{0},t]\times\mathbb{R}^{3}}Q^{\alpha\beta}\bar{\pi}_{\alpha\beta}\,d\omega+\frac{n-1}{2}\int_{[t_{0},t]\times\mathbb{R}^{3}}\psi^{2}\Box_{h}(t)\,d\omega.
\]
\end{proposition}

步骤 3. 为证明有界性定理, 需控制 $\int_{[t_0,t]\times\mathbb{R}^3}Q^{\alpha\beta}\bar{\pi}_{\alpha\beta}$. 相对零标架展开, 有
\[
Q^{\alpha\beta}\bar{\pi}_{\alpha\beta}
=\frac{1}{4}Q_{44}\bar{\pi}_{33}+\frac{1}{4}Q_{33}\bar{\pi}_{44}+\frac{1}{2}Q_{34}\bar{\pi}_{34}
-Q_{3A}\bar{\pi}_{4A}-Q_{4A}\bar{\pi}_{3A}+Q_{A B}\bar{\pi}_{A B}.
\]
我们需要形如
\begin{align*}
\int_{[t_{0},t]\times\mathbb{R}^{3}}Q^{\alpha\beta}\bar{\pi}_{\alpha\beta}
&\lesssim\lambda^{-\epsilon}\sup_{[t_{0},t]}\int_{\Sigma_{\tau}}\bigl(\underline{u}^{2}(D_{4}\psi)^{2}+u^{2}(D_{3}\psi)^{2}+\underline{u}^{2}|\nabla\psi|^{2}+\psi^{2}\bigr)\\
&=\lambda^{-\epsilon}\sup_{[t_{0},t]}\mathcal{E}_{0}[\psi](\tau)\lesssim\lambda^{-\epsilon}\sup_{[t_{0},t]}\mathcal{Q}_{0}[\psi](\tau)
\end{align*}
的估计, 对任意 $\epsilon > 0$. 特别考虑项
\[
\int_{[t_{0},t]\times\mathbb{R}^{3}}Q_{33}\bar{\pi}_{44}=\int_{[t_{0},t]\times\mathbb{R}^{3}}\left|D_{3}\psi\right|^{2}\bar{\pi}_{44}.
\]
为得到所需估计, 我们需要
\[
|\bar{\pi}_{44}(\tau,\cdot)|\lesssim u^{2}\tau^{-1}\lambda^{-\epsilon}.
\]
换言之在波区 $|u| << t$ 外 $\bar{\pi}_{44}$ 必须很小. 事实上为估计 $\int Q^{\alpha\beta}\bar{\pi}_{\alpha\beta}$ 中多数项, 我们需要对 $K$ 的修正形变张量 $\bar{\pi}$ 的如下估计.
\setcounter{mtthm}{17}
\begin{proposition}\label{prop:3.18}
$K$ 的修正形变张量的零分量满足估计:
\[
|\bar{\pi}_{44}|\leq u^{2}\lambda^{-\epsilon}\tau^{-1},\qquad|\bar{\pi}_{34}|\leq\underline{u}^{2}\lambda^{-\epsilon}\tau^{-1},
\]
\[
|\bar{\pi}_{33}|\leq\underline{u}^{2}\lambda^{-\epsilon}\tau^{-1},\quad|\bar{\pi}_{3A}|\leq\underline{u}^{2}\lambda^{-\epsilon}\tau^{-1},
\]
\[
|\bar{\pi}_{4A}|\leq u\underline{u}\lambda^{-\epsilon}\tau^{-1},\quad|\bar{\pi}_{A B}|\leq\underline{u}^{2}\lambda^{-\epsilon}\tau^{-1}.
\]
\end{proposition}

步骤 4. 上述估计可处理除项
\[
\frac{1}{2}\int_{[t_{0},t]\times\mathbb{R}^{3}}\bar{\pi}D_{3}\psi D_{4}\psi,\qquad\text{其中}\quad \bar{\pi}=\delta^{A B}\bar{\pi}_{A B}
\]
外的所有项. 为估计此项我们需要分部积分论证, 它将引入 $D_4\bar{\pi}$ 与 $\bar{\nabla}\bar{\pi}$, 对它们我们需要类似于\propref{prop:3.18} 的估计.
步骤 5. 我们需要证明\propref{prop:3.18} 的渐近估计, 以及对 $D_4\bar{\pi}$ 与 $\bar{\nabla}\pi$ 的类似估计. 这依赖于对程函方程 (79) 的关键分析. 事实上易见, 由直接计算 $\bar{\pi}$ 依赖于时间叶层 $\Sigma_t$ 的第 II 基本形式 $k$ 以及更重要的 $u$ 的 Hessian $D^2u$. 由于 $k_{ij} = -\frac{1}{2}\partial_t h_{ij}$, 我们对 $k$ 有很好控制. 另一方面 $u$ 的 Hessian 满足 Ricatti 型方程
\[
D_{L}(D^{2}u)+D^{2}u\cdot D^{2}u=R(L,L),
\]
其中 $R(L,L)$ 为曲率张量 $R_{\alpha\beta\gamma\delta}$ 与 $L$ 二次缩并的分量. 初看似乎 $D^{2}u$ 具有与曲率张量 $R$ 相同的可微性. 而且为在某时刻 $t$ 估计 $D^{2}u$, 需在零锥 $C_{u}$ 的测地线生成元上积分 $R(L,L)$. 因此, 由假设 (138), 由于 $R$ 依赖于 $h$ 的两个导数, 我们似乎能得到的最佳估计是
\[
\|D^{2}u(t)\|_{L^{\infty}}
\lesssim\int_{0}^{t}|\partial^{2}h(\tau,\cdot)|\,d\tau
\lesssim\lambda^{-2(1-a)}.
\]
沿此路线的论证只会给出 Tataru 得到的最优结果. 要做得更好我们必须利用非线性假设 (142). 能利用该假设是我们几何方法的关键优势.
为此我们需要相对典范零标架分解 $D^{2}u$. 这可通过引入如下量实现:
\[
\angle_{A B}=\langle D_{A}e_{4},e_{B}\rangle,\qquad \eta_{A}=\frac{1}{2}\langle D_{3}e_{4},e_{A}\rangle.
\]
又令
\[
b^{-1}=\partial_{t}u=|\nabla u|_{h}.
\]
我们把 $\chi_{AB}$ 分解为 $\mathrm{tr}\chi = \delta^{AB}\chi_{AB}$ 与无迹部分 $\hat{\chi}_{AB}$. 它们满足方程:
\[
L(b)=-b\, k_{N N},
\]
\[
L(\mathrm{tr}\chi)+\frac{1}{2}(\mathrm{tr}\chi)^{2}=-|\hat{\chi}|^{2}-k_{N N}\mathrm{tr}\chi-R_{44},
\]
\[
\mathcal{P}_{4}\hat{\chi}_{A B}+\frac{1}{2}(\mathrm{tr}\chi)\hat{\chi}_{A B}=-k_{N N}\hat{\chi}_{A B}-\hat{\alpha}_{A B},
\]
\[
\mathcal{P}_{4}\eta_{A}+\frac{1}{2}(\mathrm{tr}\chi)\eta_{A}=-(k_{B N}+\eta_{B})\hat{\chi}_{A B}-\frac{1}{2}(\mathrm{tr}\chi)k_{A N}-\frac{1}{2}\beta_{A},
\]
其中 $\alpha_{AB} = R_{A4B4}$, $\hat{\alpha} = \delta^{AB} \alpha_{AB}$, $\beta_{A} = R_{A43B}$. 若我们试图从各自输运方程 (89)、(90) 与 (91) 估计 $\mathrm{tr}\chi$、$\hat{\chi}$、$\eta$, 会遇到上面提到的同样困难. 我们用如下两个思想绕过这些困难:
思想 1. 注意出现在 $\mathrm{tr}\chi$ 输运方程中的曲率项 $R_{44}$ 有特殊结构. 事实上我们有

\setcounter{mtthm}{18}
\begin{lemma}[显著分解]\label{lem:3.19}
$R_{44} = e_4^\mu e_4^\nu R_{\mu\nu}$(Ricci 曲率的)分量容许如下分解:
\[
R_{44}=L(z)-\frac{1}{2}e_{4}^{\mu}e_{4}^{\nu}\Box_{h}H_{\mu\nu}+Error,
\]
其中函数 $z$ 与 Error 满足点态估计 $|z| \lesssim |\partial H|$, $|Error| \lesssim (\partial H)^2$.\end{lemma}

由该引理我们可把 $\operatorname{tr}\chi$ 的输运方程改写为:
\setcounter{mtthm}{19}
\begin{proposition}\label{prop:3.20}
令 $y = \operatorname{tr}\chi - \tfrac{2}{s}$, 其中 $s = t - u$, $z$ 与 error 如上面引理. 则
\[
L(y+z)+\operatorname{tr}\chi(y+z)
=\frac{1}{2}(y+z)^{2}+\frac{2}{s}z-\frac{1}{2}z^{2}-|\hat{\chi}|^{2}
-k_{N N}\operatorname{tr}\chi+\frac{1}{2}e_{4}^{\nu}e_{4}^{\mu}\Box_{h}H_{\mu\nu}-Error.
\]
\end{proposition}

思想 2: $y + z$ 的方程看来比 $\text{tr}\chi$ 的方程好得多. 不幸的是对 $\hat{\chi}$ 与 $\eta$ 不能做同样技巧. 而且 $y + z$ 的方程本身与 $\hat{\chi}$ 的输运方程耦合. 为解决这个问题我们求助于一个完全不同的几何方程, 即 $\chi$ 的 Codazzi 方程. 它有形式
\[
(\mathrm{div}\,\hat{\chi})_{A}+\hat{\chi}_{A B}k_{B N}=\frac{1}{2}(\nabla_{A}\mathrm{tr}\chi+k_{A N}\mathrm{tr}\chi)-R_{B4A B}.
\]
注意方程 $\mathrm{div}\,\hat{\chi}_{A}=F_{A}$ 是曲面 $S_{t,u}$ 上的椭圆 Hodge 系统. 利用上述 $y+z$ 的输运方程与 (93) 的耦合系统, 我们导出形如
\[
|\mathrm{tr}\chi-\frac{2}{s}|\lesssim\lambda^{-\bar{a}},\qquad |\hat{\chi}|\lesssim\lambda^{-\bar{a}}
\]
的估计, 其中
\[
\bar{a}=1-\frac{a^{2}}{2}.
\]
这远好于前面得到的 $\lambda^{-2(1-a)}$ 估计.
\setcounter{mtthm}{20}
\begin{remark}\label{rem:3.21}
对 $\eta$ 我们有类似方法. 思想是得到 $\eta$ 的散度--旋度系统, 耦合于一个类似于 $\text{tr}\chi$ 的输运方程.\end{remark}

\setcounter{mtthm}{21}
\begin{remark}\label{rem:3.22}
利用耦合椭圆 Hodge 系统与有利输运方程的思想, 出现在非常不同背景的 \cite{klai12} 中, 与 Minkowski 空间的全局稳定性有关.\end{remark}
